\chapter{Existing Approach for CNOT Gate Composition}

In this section, we review existing methods for CNOT circuit synthesis that aim to optimize circuit depth.

\section{Cost-Minimization Method}
De Brugière et al. \cite{9475957} addressed the problem of minimizing the depth of linear reversible quantum circuits, motivated by the observation that circuit depth is closely tied to execution time under limited qubit decoherence. Their work reformulates the synthesis of a CNOT circuit as the problem of reducing an invertible Boolean matrix to the identity through elementary row and column operations, and proposes two complementary strategies to solve it.

The first, DaCSynth, is a divide-and-conquer framework that recursively partitions the matrix into blocks and eliminates them using a bipartite-graph matching approach, yielding a provable worst-case depth upper bound of $\frac{4}{3}n + 8\lceil \log_2(n) \rceil$ for an $n$-qubit operator, an improvement over prior asymptotic and practical bounds for intermediate-sized registers. The second is a family of purely greedy, cost-minimization algorithms that iteratively apply row and column operations to reduce a chosen cost function toward its minimum, reached when the matrix becomes a permutation matrix. Four cost functions are considered, based on the sum and logarithmic Hamming weight of matrix rows (and their inverses), which balance encouraging overall sparsification against prioritizing "nearly complete" rows to avoid local minima in Table \ref{tab:cost-function}.

\begin{table}
    \centering
    \begin{tabularx}{\textwidth}{C|C}
        \hline
        Cost Function  & Linear Reversible Cost Function\\
         \hline
         $$h_{sum}(A) = \sum_{i,j}a_{i,j}$$ & $$H_{sum}(A) = h_{sum}(A) + h_{sum}(A^{-1})$$\\
         \hline
         $$h_{prod}(A) = \sum_{i=1}^{n}log_{2}(\sum_{j=1}^{n}a_{i,j})$$ & $$H_{prod}(A) = h_{prod}(A) + h_{prod}(A^{-1})$$\\
    \hline
\end{tabularx}
    \caption{Cost function in \cite{10.1145/3474226}}
    \label{tab:cost-function}
\end{table}

Their benchmarks show that while DaCSynth performs consistently well across problem sizes and provides strong worst-case guarantees, the greedy methods outperform DaCSynth for small or structurally simple operators, though their performance degrades sharply as matrix size grows due to the difficulty of escaping local minima in the cost landscape. Applied to a library of reversible functions (including Galois field multipliers and arithmetic circuits), their combined framework achieves depth reductions of up to 92\% (ancilla-free) and 99\% (with ancillary qubits) relative to prior state-of-the-art methods.

This greedy cost-minimization framework forms the direct algorithmic foundation for the depth-oriented approach adopted in this thesis: subsequent work by Shi and Feng \cite{cryptoeprint:2024/381}, and our own enhanced algorithm, build on and extend the cost functions and search strategy first introduced here.

\section{Depth-Oriented Method}

\begin{table}[htbp]
\centering
\begin{tabular}{c|c}
\hline
Cost Function  & Linear Reversible Cost Function\\
\hline
\multirow{2}{*}{$h_{sq}(A) = \sum_{i}(\sum_{j}a_{i,j})^{2}$} & $H_{sqr}(A) = h_{sq}(A) + h_{sq}((A^{-1})^{T})$ \\ 
\cline{2-2}
 & $H_{sqc}(A) = h_{sq}(A^{T}) + h_{sq}(A^{-1})$ \\ 
\hline
\end{tabular}
	\caption{Square every row's Hamming-weight to enlarge uncomplete part.}
	\label{tab:square-hamming-weight}
\end{table}

Shi and Feng \cite{cryptoeprint:2024/381} extended the depth-oriented greedy framework of de Brugière et al. to specifically target the synthesis of low-depth CNOT circuits for symmetric-cipher linear layers, with a particular focus on AES MixColumns. Their central algorithmic contribution refines the original greedy cost-minimization strategy in three respects: (1) in addition to the logarithmic Hamming-weight cost function, they introduce a cost function based on the square of each row's Hamming weight, which more aggressively prioritizes rows that are "far from complete" in Table \ref{tab:square-hamming-weight}; (2) row and column operations are evaluated with distinct cost functions, so that column-wise sparsity is explicitly accounted for during column operations rather than only through the logic used for rows; and (3) they derive a necessary and sufficient condition for determining whether a matrix can be implemented with depth exactly 1, allowing sparse matrices to be handled more efficiently in the final stage of the search~\eqref{eq:depth-one-cost-function}.

\begin{equation}
	\label{eq:depth-one-cost-function}
	\begin{split}
	H_{prodr}(A) &= h_{prod}(A) + h_{prod}((A^{-1})^{T}) \\
	H_{prodc}(A) &= h_{prod}(A^{T}) + h_{prod}(A^{-1})
	\end{split}
\end{equation}

Applying this improved algorithm to AES MixColumns, they report a CNOT circuit of depth 10, improving on the previous best result of depth 16 reported by Liu et al.\cite{cryptoeprint:2023/1417}, and also outperforming a direct application of de Brugière et al.'s original greedy algorithm, which yields depth 12 on the same matrix. Beyond AES, they apply their method to a broad collection of MDS matrices and linear layers used in other block ciphers (including ANUBIS\cite{barreto2000anubis}, CLEFIA, TWOFISH, and several lightweight ciphers), consistently achieving the lowest reported depths across nearly all cases. In addition to circuit-depth optimization, the paper proposes a compressed pipeline structure for synthesizing full AES encryption circuits and Grover/Encryption oracles, which reduces the number of intermediate qubit registers relative to the standard pipeline structure while retaining comparable round depth, and reports the lowest published T-depth for an AES-128 Encryption oracle to date.

We clarify their work that as pseudocode as Algorithm~\ref{alg:Modified-Shi-Feng-Greedy}

\begin{algorithm}
\caption{Row/Col Greedy}
\label{alg:Modified-Shi-Feng-Greedy}
\begin{algorithmic}[1]
\Require $M \in \mathbb{F}_{2}^{n \times n}$, $M' \in \mathbb{F}_{2}^{n \times n}$, $L_{r}, L_{c}, Ls_{r}, Ls_{c}, p$ and over\_depth
\Ensure $Ls_{r}, Ls_{c}, row_{op}, col_{op}$ that implement permutation matrix with length $d$
\State SIZE $\leftarrow$ length($M$)
\State minm $\leftarrow$ system.float\_info.max
\State select\_list $\leftarrow$ list()
\State $visi_{row}$ $\leftarrow [\varnothing]$, $visi_{col}$ $\leftarrow [\varnothing]$
\State close\_permu $\leftarrow$ False
\State LIMIT $\leftarrow$ LATEST\_MINM\_DEPTH
\While{$M$ != Permutation}
	\State select\_list.clear()
	\State $L_{row}$.clear(), $L_{col}$.clear(), $L_{row}^{cost}$.clear(), $L_{col}^{cost}$.clear()
	
	\State $H_{r}, H_{c}$, minm $\leftarrow$ functionCost($M, M', p$)
    
    \If{!close\_permu}
		\State $L_{col} \leftarrow L\_Collection(visi_{row})$
	    \State select\_list, minm $\leftarrow Row\_Op\_Sel( L_{col}, L_{col}^{cost}, M, M', p)$ 
    \EndIf
	\State $L_{row} \leftarrow L\_Collection(visi_{row}$)
	\Function{availableRowCollection}{$visi_{row}$, SIZE}
    \State \Return $L_{row}$
	\EndFunction
	\Function{availableRowOperatorSelection}{$L_{row}$, $L_{row}^{cost}$, $mat$, $inv$, $p$}
    \State \Return select\_list, minm
	\EndFunction
    
	\Function{availableOperatorExecution}{select\_list, $L_{r}, L_{c}, Ls_{r}, Ls_{c}$, $mat$, $inv$, $row_{op}, col_{op}$, $visi_{row}, visi_{col}$, $d$, SIZE, one}
    \State \Return $L_{r}, L_{c}, Ls_{r}, Ls_{c}$, $mat$, $inv$, $row_{op}, col_{op}$, $visi_{row}, visi_{col}$, $d$, one
	\EndFunction
    
    \If{$d >$  LIMIT}
	    \State $flag$ $\leftarrow$ True
	    \State \Return
    \EndIf
\EndWhile
\State \Return $L_{r}, L_{c}, Ls_{r}, Ls_{c}, mat, row_{op}, col_{op}, d$, flag
\end{algorithmic}
\end{algorithm}